\section{总结与延伸}

\subsection{核心结论}

最后这一节把全篇从产业叙事重新压缩成可复用判断。读者不需要记住每个公司数字，但应该带走一套判断 AI 竞争的层次图：资本开支、模型范式、入口、数据和创业生态怎样互相约束。

本期的主线可以压缩成九个结论。

\begin{enumerate}
\item AI Bubble 只能解释估值情绪，AI War 才能解释巨头和国家为什么持续投入。
\item OpenAI 的看得见收入来自订阅、广告电商和 API，看不见收入来自 Agent 替代劳动和 AI 做高价值增量任务。
\item Nvidia GPU 与 Google TPU 代表开放生态和垂直整合两条硬件/云/模型路线。
\item GPT、Claude、Gemini 的交替领先会成为常态，直到下一范式真正拉开差距。
\item 基础模型公司像综合电商，Scale Data 的效率很高，创业公司必须找到差异化数据和工作流。
\item Google 王者归来不等于危机解除，ChatGPT 也不只是聊天工具；二者都在争夺搜索和个人助理入口。
\item Pre-training 像石油，RL 专家数据像新能源，Online Learning 像尚未突破的核聚变。
\item “模型即产品，数据即模型”意味着 AI 产品的护城河要回到数据分布、用户反馈和专家过程。
\item 华人 AI 创业者的机会来自全球市场、中国人才、工程能力和产品说话，而不是身份包装。
\end{enumerate}

\subsection{开放问题}

这些问题决定 2026 年之后的分叉。

\begin{itemize}
\item Online Learning 会在 2026 年出现清晰信号，还是继续停留在研究叙事？
\item ChatGPT 融合搜索的速度，会不会快于 Google 自我革新的速度？
\item Gemini 的多模态和 Android 生态，能否转化为高频个人助理心智？
\item Anthropic 的 Coding 和企业战略 Beta，能否继续抵抗 OpenAI 和 Google 的横向扩张？
\item Agentic Web 的支付、授权、人机验证和网站协议，谁会率先标准化？
\item Robotics 的数据路线会收敛，还是继续长期分化？
\item 中国团队能否在全球 AI 应用独角兽中重新获得互联网时代那样的占比？
\end{itemize}

\subsection{拓展阅读}

\begin{itemize}
\item EP136 全球大模型季报第 9 集：Coding、模型 OS 与硅谷御三家。
\item EP134 数据综述：数据金字塔、Recipe、标注与定价。
\item EP139 Agent 技术史综述：Agent 的技术谱系和社会辐射。
\item EP130 张月光访谈：AI Native 产品、上下文设计和 Agent 产品方法。
\item OpenAI、Anthropic、Google DeepMind、SSI、Thinking Machines 相关公开访谈与技术报告：理解下一范式和 Post-training 路线。
\end{itemize}

\begin{importantbox}{最后的判断}
EP127 最值得留下的是一个跨层框架：AI 的竞争不是单纯模型榜单，也不是单纯资本泡沫，而是算力、数据、入口、产品、学习范式和国家战略同时变化。未来赢家不一定是某个季度 benchmark 最强的公司，而是能把真实数据、持续学习和用户入口连成闭环的系统。
\end{importantbox}

\end{document}
